\documentclass[11pt]{amsart}
\usepackage[T1]{fontenc}
\usepackage{amsmath,amssymb,amsthm}
\usepackage{graphicx}
\usepackage[colorlinks=true,linkcolor=blue,citecolor=blue,urlcolor=blue]{hyperref}

\newtheorem{theorem}{Theorem}
\newtheorem{lemma}{Lemma}
\newcommand{\R}{\mathbb R}
\newcommand{\gammaRho}{\gamma_{n,\rho}}

\begin{document}

\title{Gaussian Multi-Set Concentration from Isoperimetric Contraction}
\author{Run fa\_banach\_001}
\date{August 11, 2026}
\maketitle
\pagestyle{plain}

\section*{Status}

This is a candidate full answer, subject to human review, to the Gaussian
multi-set concentration conjecture in Section 5.1 of Gozlan and Herry
\cite{GH}.  It gives the requested profile with the stronger constant
$C_{k,\rho,n}=\rho/2$, independent of $k$ and $n$.

\section*{Source Question}

For $k\geq1$, let $\Delta_k$ consist of all
$a=(a_1,\ldots,a_k)\in[0,1]^k$ such that
\[
 \sum_{j=1}^k a_j\leq1,
 \qquad
 a_i+\sum_{j=1}^k a_j\geq1
 \quad(1\leq i\leq k).
\]
A probability metric space has the multi-set concentration property of order
$k$ with profile $\alpha_k$ if, whenever the Borel sets
$A_1,\ldots,A_k$ are pairwise positively separated and
$(\mu(A_1),\ldots,\mu(A_k))\in\Delta_k$, their union $A$ satisfies
\begin{equation}\label{eq:multiset}
 \mu(A_r)\geq
 1-(1-\mu(A))\alpha_k(r),
 \qquad
 0<r\leq\frac12\min_{i\ne j}d(A_i,A_j).
\end{equation}
Here $A_r=\{x:d(x,A)<r\}$.

The source asks whether every probability measure
$d\mu=e^{-V}dx$ on $\R^n$ with $V$ smooth and
$\operatorname{Hess}V\geq\rho I>0$ has this property with a Gaussian
profile $\exp(-C_{k,\rho,n}r^2)$.

\begin{center}
\includegraphics[width=\textwidth]{figures/open_problem_crop.png}
\end{center}

\section*{Proof Intuition}

The condition defining $\Delta_k$ already forces the union of the $k$ sets
to have mass at least $k/(k+1)$, and hence at least $1/2$.  Strong
log-concavity makes $\mu$ a $1$-Lipschitz image of the Gaussian measure
with covariance $\rho^{-1}I$ by Caffarelli's contraction theorem.  Gaussian
isoperimetry therefore bounds the enlargement of the \emph{union}; no
separate higher-order estimate is needed.  A one-line shifted-tail
calculation supplies exactly the multiplicative factor $1-\mu(A)$ required
in \eqref{eq:multiset}.

\section*{Two Elementary Lemmas}

Let $\Phi$ and $\phi$ denote the distribution and density functions of a
standard one-dimensional Gaussian, and let
$\overline\Phi=1-\Phi$.

\begin{lemma}[The $\Delta_k$ mass threshold]\label{lem:mass}
If $a=(a_1,\ldots,a_k)\in\Delta_k$ and
$s=\sum_{i=1}^k a_i$, then
\[
 s\geq\frac{k}{k+1}\geq\frac12.
\]
\end{lemma}

\begin{proof}
Put $b=1-s$.  The defining inequalities give $a_i\geq b$ for every $i$.
Thus $s\geq kb=k(1-s)$, which is equivalent to the claim.
\end{proof}

\begin{lemma}[Multiplicative Gaussian tail shift]\label{lem:tail}
For $z,t\geq0$,
\[
 \overline\Phi(z+t)
 \leq e^{-t^2/2}\overline\Phi(z).
\]
\end{lemma}

\begin{proof}
Changing variables $u=v+t$ gives
\[
 \overline\Phi(z+t)
 =\int_z^\infty\phi(v+t)\,dv
 =e^{-t^2/2}\int_z^\infty e^{-tv}\phi(v)\,dv
 \leq e^{-t^2/2}\overline\Phi(z),
\]
because $v\geq z\geq0$.
\end{proof}

\section*{Full Answer}

\begin{theorem}\label{thm:answer}
Let $\mu$ be a probability measure on $\R^n$ with smooth density
$e^{-V}$ and suppose
\[
 \operatorname{Hess}V\geq\rho I
\]
for some $\rho>0$.  For every $k\geq1$, the space
$(\R^n,|\cdot|,\mu)$ has the multi-set concentration property of order
$k$ with
\[
 \alpha_k(r)=\exp\!\left(-\frac{\rho r^2}{2}\right),
 \qquad r\geq0.
\]
In particular, the conjecture of \cite[Section 5.1]{GH} holds with
$C_{k,\rho,n}=\rho/2$.
\end{theorem}

\begin{proof}
Let $\gammaRho$ be the centered Gaussian probability measure on $\R^n$
with covariance $\rho^{-1}I$.  Since
\[
 V(x)-\frac\rho2|x|^2
\]
is convex, Caffarelli's contraction theorem \cite{Caffarelli} supplies a
$1$-Lipschitz transport map $F$ with
$F_\#\gammaRho=\mu$.

We first record the resulting Gaussian isoperimetric comparison.  For a
Borel set $E\subset\R^n$, put $B=F^{-1}(E)$.  The Lipschitz property gives
$B_r\subset F^{-1}(E_r)$, so Gaussian isoperimetry yields
\begin{equation}\label{eq:comparison}
 \mu(E_r)
 \geq \gammaRho(B_r)
 \geq
 \Phi\!\left(\Phi^{-1}(\mu(E))+\sqrt\rho\,r\right).
\end{equation}

Now take positively separated Borel sets $A_1,\ldots,A_k$ satisfying the
hypotheses of the multi-set property and set $A=\bigcup_iA_i$.  Positive
separation makes the sets disjoint, hence, with
$a_i=\mu(A_i)$,
\[
 s:=\mu(A)=\sum_{i=1}^k a_i.
\]
Lemma \ref{lem:mass} gives $s\geq1/2$.  Therefore
$z=\Phi^{-1}(s)\geq0$.  Taking complements in
\eqref{eq:comparison}, and applying Lemma \ref{lem:tail} with
$t=\sqrt\rho\,r$, gives
\begin{align*}
 1-\mu(A_r)
 &\leq \overline\Phi(z+\sqrt\rho\,r)\\
 &\leq e^{-\rho r^2/2}\overline\Phi(z)\\
 &=(1-\mu(A))e^{-\rho r^2/2}.
\end{align*}
This is exactly \eqref{eq:multiset} with the asserted profile.  In fact the
bound holds for every $r\geq0$, stronger than the separation-limited range
required by the definition.
\end{proof}

\section*{Verification Notes}

The proof uses only two standard sharp inputs: Caffarelli contraction from
$\gammaRho$ to a measure whose potential is at least as convex, and Gaussian
isoperimetry.  All remaining steps are displayed.  The direction of the
neighborhood inclusion is important: if $x$ is within $r$ of
$F^{-1}(E)$, then $F(x)$ is within $r$ of $E$, hence
$(F^{-1}E)_r\subset F^{-1}(E_r)$.

The multiplicative prefactor is not discarded.  Since
$\overline\Phi(\Phi^{-1}(s))=1-s$, Lemma \ref{lem:tail} produces precisely
the complement mass in the source definition.  The assumption
$s\geq1/2$ is exactly what ensures $z\geq0$ in that lemma.

For a half-space under $\gammaRho$ with mass $1/2$, the exponent
$\rho/2$ is asymptotically sharp as $r\to\infty$, so no larger uniform
quadratic coefficient can hold in this formulation without an additional
prefactor or linear term.

\section*{Novelty Check and Scope}

A bounded search was performed on August 11, 2026.  The local run indexes,
solution registry, attempts, and parsed-source scan were searched for
arXiv:1804.06133 and combinations of ``Gaussian multi-set concentration'',
``Caffarelli contraction'', and ``Gaussian isoperimetry''.  External
exact-phrase searches found the arXiv and 2020 Potential Analysis versions,
the authors' publication pages, a thesis reproducing the same open question,
and later work on adjacent higher-order eigenvalue bounds, but no paper
claiming this answer or applying the contraction argument to the exact
multi-set profile.  Novelty is plausible, not certified; specialist review
of MathSciNet and zbMATH remains advisable.

The theorem answers the first question in Section 5.1.  It does not answer
the separate Section 5.2 converse problem relating multi-set concentration
to lower bounds on higher eigenvalues, nor does it construct new modified
functional inequalities.

\section*{Human Review Recommendation}

Check the convention and scaling in the cited form of Caffarelli's
contraction theorem, then verify the neighborhood inclusion
$(F^{-1}A)_r\subset F^{-1}(A_r)$.  The rest reduces to Lemmas
\ref{lem:mass} and \ref{lem:tail}.  If these checks pass, Theorem
\ref{thm:answer} is a complete and stronger-than-requested solution.

\begin{thebibliography}{9}

\bibitem{Caffarelli}
L. A. Caffarelli,
\emph{Monotonicity properties of optimal transportation and the FKG and
related inequalities},
Comm. Math. Phys. \textbf{214} (2000), 547--563.

\bibitem{GH}
N. Gozlan and R. Herry,
\emph{Multiple sets exponential concentration and higher order eigenvalues},
Potential Anal. \textbf{52} (2020), 203--221;
arXiv:1804.06133.

\bibitem{Ledoux}
M. Ledoux,
\emph{The Concentration of Measure Phenomenon},
Mathematical Surveys and Monographs 89, American Mathematical Society, 2001.

\end{thebibliography}

\end{document}
